%=====================================================================
\chapter{Lightweight and Secure Virtualization}\label{ch:lightweight}
%=====================================================================
\locator{3}{Part III --- Containers and Lightweight Virtualization}

\begin{outcomes}
By the end of this chapter you will be able to:
\begin{tight}
  \item \textbf{State} the isolation--density--start-up trade and \textbf{place}
        any runtime on it.
  \item \textbf{Explain} how a microVM reaches a 125 ms boot, and \textbf{name}
        what was removed to get there.
  \item \textbf{Contrast} gVisor's approach --- intercepting system calls in
        user space --- with Kata's, which puts a kernel underneath each
        container.
  \item \textbf{Describe} a unikernel and \textbf{say} why the idea has not
        displaced the alternatives.
  \item \textbf{Choose} a runtime for a stated workload and \textbf{quantify}
        what the choice costs.
\end{tight}
\end{outcomes}

\begin{prereq}
\chref{containers} (the shared kernel and the 350-call attack surface) and
\chref{platforms} (what a hypervisor contains, and Xen's disaggregation
argument). \chref{paravirt} is also assumed: every design here is fast because
it has virtio and nothing else.
\end{prereq}

\begin{keyterms}
microVM $\bullet$ Firecracker $\bullet$ Cloud Hypervisor $\bullet$ device model
$\bullet$ boot time $\bullet$ jailer $\bullet$ gVisor $\bullet$ Sentry $\bullet$
Gofer $\bullet$ system-call interception $\bullet$ Kata Containers $\bullet$
RuntimeClass $\bullet$ unikernel $\bullet$ library operating system $\bullet$
WebAssembly $\bullet$ WASI $\bullet$ attack surface $\bullet$ cold start
$\bullet$ oversubscription
\end{keyterms}

\section{The Choice Nobody Wanted to Make}

Part~II gave an isolation boundary enforced by the processor, at a few hundred
megabytes and half a minute per instance. Part~III gave a boundary enforced by
the kernel, at a few megabytes and a tenth of a second.

For most workloads the second is right. For one case it is unacceptable:
running \emph{other people's} code. A public cloud function service executes
whatever a stranger uploads, beside another stranger's, on hardware shared with
a third. The 350-call Linux system-call interface is then not an abstraction
--- it is the perimeter, and it is enormous.

The honest statement of the problem is that this is a three-way trade and you
may normally have two.

\dsfig{%
\begin{tikzpicture}[font=\scriptsize,
  ax/.style={neutral!60, line width=0.9pt},
  pt/.style={circle, inner sep=1.6pt, line width=0.8pt},
  lb/.style={font=\scriptsize\bfseries, align=center},
  note/.style={font=\scriptsize\itshape, align=center}
]
  \draw[ax, -{Stealth[length=2.2mm]}] (0,0) -- (12.2,0);
  \draw[ax, -{Stealth[length=2.2mm]}] (0,0) -- (0,5.0);
  \node[font=\scriptsize\bfseries, text=neutral!40!black, anchor=north]
       at (6.1,-0.35) {start-up time and overhead per instance $\longrightarrow$};
  \node[font=\scriptsize\bfseries, text=neutral!40!black, rotate=90,
        anchor=south] at (-0.45,2.5) {strength of the isolation boundary
        $\longrightarrow$};

  \foreach \x/\t in {1.3/{100 ms\\4 MB}, 4.2/{125 ms\\5 MB},
                     7.0/{250 ms\\50 MB}, 10.3/{25 s\\400 MB}}{
    \node[note, text=neutral, anchor=north] at (\x,-0.05) {\t};}

  % points
  \node[pt, draw=greenhl, fill=greenhl!30] (c)  at (1.3,1.0) {};
  \node[pt, draw=goldhl,  fill=goldhl!30]  (g)  at (7.0,2.6) {};
  \node[pt, draw=purplehl,fill=purplehl!30](f)  at (4.2,3.7) {};
  \node[pt, draw=purplehl,fill=purplehl!30](k)  at (7.6,3.9) {};
  \node[pt, draw=secondary,fill=secondary!30](v) at (10.3,4.3) {};
  \node[pt, draw=accent,  fill=accent!30]  (u)  at (2.4,3.2) {};

  \node[lb, text=greenhl!45!black, anchor=west]  at (1.55,0.75)
       {container\\{\scriptsize\mdseries shared kernel}};
  \node[lb, text=goldhl!45!black, anchor=west]   at (7.25,2.35)
       {gVisor\\{\scriptsize\mdseries kernel in user space}};
  \node[lb, text=purplehl!70!black, anchor=east] at (3.95,3.95)
       {Firecracker\\{\scriptsize\mdseries microVM}};
  \node[lb, text=purplehl!70!black, anchor=west] at (7.85,4.15)
       {Kata\\{\scriptsize\mdseries VM per container}};
  \node[lb, text=secondary!80!black, anchor=east] at (10.1,4.6)
       {full VM};
  \node[lb, text=accent!70!black, anchor=west]   at (2.65,3.0)
       {unikernel\\{\scriptsize\mdseries one app, no kernel}};

  \draw[greenhl!50, line width=0.9pt, dashed] (1.3,1.0) -- (4.2,3.7);
  \node[note, text=accent, anchor=west] at (4.4,1.6)
       {Part III's whole answer:\\
        move up and left --- keep the\\
        hardware boundary, pay the\\
        container's start-up cost};
\end{tikzpicture}}%
{The trade, with measured figures on the horizontal axis. The dashed line is
the move the designs in this chapter make: a hardware-enforced boundary at
something close to a container's cost.}{isolation-density}

\section{MicroVMs}

A virtual machine is slow to start and heavy because of what the hypervisor
emulates: a BIOS, a PCI bus, an IDE controller, a VGA adapter, a PS/2 keyboard
--- a 1990s PC, assembled in software so that an unmodified operating system
recognises it. \chref{paravirt} showed this is also where the throughput goes.

A \term{microVM} deletes all of it. If the guest is going to use virtio for
everything anyway, the emulated chipset is dead weight.

\begin{definitionbox}[MicroVM]
A \term{microVM} is a virtual machine with a minimal device model: typically
virtio-net, virtio-block, a serial console and a one-shot timer, and nothing
else. No BIOS, no PCI enumeration, no legacy devices. The kernel is loaded
directly rather than booted through firmware.

\smallskip
\term{Firecracker}, written at Amazon in Rust and released in 2018, is the
reference example. It is about 50{,}000 lines against QEMU's 1.4 million, boots
a Linux guest in roughly \textbf{125 ms}, and adds under \textbf{5 MB} of
memory per instance. It runs AWS Lambda and AWS Fargate: every function
invocation you have ever made on Lambda ran in one.
\end{definitionbox}

\begin{conceptbox}[Where the Boot Time Went]
A conventional guest spends its 25 seconds roughly like this: firmware
initialisation and PCI enumeration, bootloader, kernel device probing for
hardware that is not there, then an init system bringing up services. Almost
none of it is the application.

\smallskip
Firecracker removes each in turn. There is no firmware, so nothing to
initialise. There is no PCI bus, so nothing to enumerate. The kernel is loaded
into memory by the VMM and entered directly, with a command line that names the
devices, so probing finds exactly what exists and stops. What remains is
roughly 125 ms, most of it the Linux kernel's own startup.

\smallskip
The lesson generalises past this chapter: \emph{the cost was compatibility, not
virtualization}. A guest willing to give up looking like a PC can have a
hardware isolation boundary for about the price of a container --- and that one
observation is what made serverless computing economic.
\end{conceptbox}

\begin{alertbox}[The jailer, and defence in depth]
Firecracker's small size is a security argument only if you also assume the VMM
process itself can be attacked. Amazon assumes it can.

\smallskip
Each Firecracker process is therefore started by a \term{jailer} that places it
in its own namespaces and cgroup, chroots it, drops every capability it does
not need and applies a seccomp filter restricting it to a few dozen system
calls. An attacker who escapes the guest lands in a process that can barely
talk to the kernel --- the container mechanisms of \chref{containers} used to
contain the hypervisor rather than the workload.

\smallskip
This is the strongest practical illustration in this book that the layers
compose. It is not virtual machines \emph{or} containers; it is a microVM
inside a container, and the combination is stronger than either.
\end{alertbox}

\section{gVisor: a Kernel in User Space}

Google's answer inverts the problem. Instead of putting a real kernel
underneath the container, \term{gVisor} puts a \emph{fake} one in front of it.

A gVisor sandbox runs the application's system calls into \term{Sentry}, a
user-space process written in Go that implements most of the Linux system-call
interface itself. Sentry handles the call in user space wherever it can; where
it genuinely needs the host, it makes a far narrower call of its own. File
access goes through a second process, the \term{Gofer}, so Sentry itself need
not hold file descriptors.

\begin{conceptbox}[What gVisor Buys and What It Costs]
\textbf{Buys:} the host kernel sees perhaps 60 distinct system calls from
Sentry rather than 350 arbitrary ones from untrusted code. The attack surface
shrinks by roughly a factor of six, without a hypervisor and without a second
kernel's memory.

\smallskip
\textbf{Costs:} every system call is now handled by a Go program instead of by
the kernel. System-call-heavy workloads --- anything doing many small reads and
writes --- lose substantially, commonly 10--30\% and occasionally much more.
CPU-bound work is barely affected, because it makes few calls.

\smallskip
\textbf{And a third thing:} Sentry implements \emph{most} of Linux, not all of
it. Unusual system calls, exotic \texttt{ioctl}s and some \texttt{/proc}
behaviour are missing or approximated, so an application can fail in ways it
never would natively. It is used in production at Google and in Google Cloud
Run; it needs testing with your workload rather than faith.
\end{conceptbox}

\section{Kata Containers}

Kata takes the most direct route: make every container a virtual machine, and
hide that fact completely.

Kata implements the OCI runtime interface, so Kubernetes drives it exactly as
it drives runc. Underneath, each Pod gets its own lightweight VM --- using
Firecracker, Cloud Hypervisor or QEMU --- with a minimal guest kernel and an
agent that starts the container inside. The user writes ordinary manifests; a
\term{RuntimeClass} selects Kata for the workloads that need it.

\begin{conceptbox}[Per-Workload Isolation Is the Deployable Idea]
The important property is not Kata's architecture but its granularity. A
cluster does not have to choose one isolation level for everything.

\smallskip
Internal services, whose code you wrote and review, run on runc at full density.
A tenant-supplied plugin, a customer's build job, an untrusted PDF converter
--- these get \texttt{runtimeClassName: kata} and a hardware boundary, at
perhaps 50 MB and 250 ms each. The manifests are otherwise identical, and the
decision is made per workload by the people who know the trust level.

\smallskip
That is the shape of the answer this chapter has been building towards:
isolation stops being an architectural commitment and becomes a per-workload
setting.
\end{conceptbox}

\section{Unikernels and WebAssembly}

Two more positions on the same axis, both further out than anything above.

A \term{unikernel} compiles the application together with just the operating
system functions it actually calls, producing a single-address-space image that
boots directly on a hypervisor. There is no shell, no package manager, no
second process and often no notion of a user. Images are measured in megabytes
and boot in milliseconds, and the attack surface is tiny because almost nothing
is there.

They have not displaced anything, and the reasons are instructive: you cannot
debug one with the usual tools, you cannot run arbitrary existing software in
one without porting it, and the build toolchains have stayed difficult. The
idea is sound and the operational cost has never been paid down.

\term{WebAssembly} outside the browser --- Wasm with \term{WASI} --- is the
current attempt at the same goal from the other direction: a sandboxed
bytecode with a deny-by-default capability model, millisecond start-up and
images in the low megabytes. It runs under Kubernetes today through
\texttt{runwasi} and similar shims. Its limitations are the obvious ones: your
application must compile to it, and the ecosystem is much younger than the
alternatives in this chapter.

\begin{worked}[Sizing a Function Service]
A provider will run untrusted functions. Expected load: 2{,}000 concurrent
instances, average lifetime 400 ms, memory limit 128 MB each, on hosts of 64
cores and 256 GB. Compare plain containers, Kata and Firecracker.

\smallskip
\textbf{Step 1 --- memory per instance.}
\begin{tight}
  \item container: $128 + 4 = 132$ MB
  \item Kata on QEMU: $128 + 50 = 178$ MB
  \item Firecracker: $128 + 5 = 133$ MB
\end{tight}

\smallskip
\textbf{Step 2 --- instances per host}, reserving 32 GB, so 224 GB usable:
\[ \text{container} = \frac{224 \times 1024}{132} = 1{,}737 \qquad
   \text{Kata} = \frac{224 \times 1024}{178} = 1{,}288 \]
\[ \text{Firecracker} = \frac{224 \times 1024}{133} = \mathbf{1{,}724} \]

\smallskip
\textbf{Step 3 --- hosts needed for 2{,}000 concurrent.}
\[ \text{container} = \left\lceil \frac{2000}{1737} \right\rceil = 2 \qquad
   \text{Kata} = \left\lceil \frac{2000}{1288} \right\rceil = 2 \qquad
   \text{Firecracker} = 2 \]
All three need two hosts. On this workload the density difference does not
change the bill at all.

\smallskip
\textbf{Step 4 --- the cost that does differ: cold start.} With a 400 ms
lifetime, start-up is not overhead, it is a large fraction of the work:
\[ \text{container} = \frac{0.10}{0.40 + 0.10} = 20\% \qquad
   \text{Kata} = \frac{0.25}{0.40 + 0.25} = 38\% \]
\[ \text{Firecracker} = \frac{0.125}{0.40 + 0.125} = \mathbf{24\%} \]
Kata spends 38\% of every host-second starting machines. Firecracker spends
24\%, nearly matching the container, and gives a hardware boundary for it.

\smallskip
\textbf{Step 5 --- the decision.} \textbf{Firecracker.} It is within 1\% of the
container's density and within 4 percentage points of its start-up overhead,
and it is the only one of the three that puts the processor between one
stranger's code and the next.

\smallskip
\textbf{Step 6 --- and the honest caveat.} These figures assume a cold start
every time. Real function services keep warm instances pooled, which cuts the
effective cold-start rate by an order of magnitude and moves the whole
comparison. Reproduce Step 4 with your measured cold-start \emph{rate}, not the
worst case, before committing hardware.
\end{worked}

\begin{pitfall}
\begin{tight}
  \item \textbf{Treating this as a ranking.} There is no best runtime. A
        microVM for an internal service is wasted money; runc for a stranger's
        code is negligence.
  \item \textbf{Assuming gVisor is free.} The system-call path is slower, and
        how much depends entirely on how many calls your workload makes.
        Measure it, do not assume 10\%.
  \item \textbf{Expecting every program to run under gVisor.} It implements
        most of Linux, not all. Test the actual application.
  \item \textbf{Thinking a microVM removes the need to harden the container
        around it.} Firecracker ships a jailer precisely because its designers
        do not make that assumption.
  \item \textbf{Choosing a unikernel without pricing the debugging.} The
        runtime properties are excellent and the day the thing misbehaves in
        production you have no shell, no \texttt{strace} and no package to
        install one from.
  \item \textbf{Applying one isolation level to a whole cluster.} RuntimeClass
        exists so that the decision can be made per workload. Most clusters
        should mix.
\end{tight}
\end{pitfall}

\begin{lab}[Booting a MicroVM, and Timing It]
Starts a Firecracker microVM and compares its boot time and memory footprint
with a full QEMU guest and a container. Needs KVM. Expect forty minutes.

\begin{lstlisting}[language=bash]
#!/usr/bin/env bash
# ch11_microvm.sh -- 25 s, 125 ms and 100 ms, measured on one machine.
set -euo pipefail
cd ~ && mkdir -p fc && cd fc

# --- 1. Firecracker needs KVM and nothing else. ------------------------
[ -w /dev/kvm ] || { echo "no /dev/kvm -- see Chapter 4 step 1"; exit 1; }
ARCH=$(uname -m)
curl -sL -o firecracker \
  "https://github.com/firecracker-microvm/firecracker/releases/latest/\
download/firecracker-${ARCH}.tgz" || true
# If that fails, the release page lists a direct tarball; extract the
# firecracker binary beside this script.
chmod +x firecracker 2>/dev/null || true

# --- 2. A kernel and a root filesystem, loaded directly. --------------
# No bootloader and no firmware: the VMM places the kernel in memory.
curl -sL -o vmlinux \
  "https://s3.amazonaws.com/spec.ccfc.min/img/quickstart_guide/\
${ARCH}/kernels/vmlinux.bin"
curl -sL -o rootfs.ext4 \
  "https://s3.amazonaws.com/spec.ccfc.min/img/quickstart_guide/\
${ARCH}/rootfs/bionic.rootfs.ext4"

# --- 3. Describe the whole machine. Note how little there is. ---------
cat > vm.json <<'JSON'
{
  "boot-source": {
    "kernel_image_path": "vmlinux",
    "boot_args": "console=ttyS0 reboot=k panic=1 pci=off"
  },
  "drives": [{
    "drive_id": "rootfs", "path_on_host": "rootfs.ext4",
    "is_root_device": true, "is_read_only": false
  }],
  "machine-config": { "vcpu_count": 1, "mem_size_mib": 128 }
}
JSON
# pci=off is the chapter: there is no PCI bus to enumerate.

# --- 4. Boot it, and time to the login prompt. ------------------------
echo "=== firecracker ==="
time timeout 20 ./firecracker --no-api --config-file vm.json \
  2>&1 | grep -m1 -E 'login:|Welcome' || true

# --- 5. The same guest under full QEMU, for comparison. ---------------
echo "=== full qemu guest ==="
time timeout 90 sudo virsh start vss1 --console 2>&1 \
  | grep -m1 'login:' || true
sudo virsh destroy vss1 || true

# --- 6. And a container, for the other end of the axis. ---------------
echo "=== container ==="
time docker run --rm alpine:3.20 /bin/true

# --- 7. Memory footprint of the VMM process itself. -------------------
./firecracker --no-api --config-file vm.json &
FCPID=$!
sleep 3
grep -E '^(VmRSS|VmSize)' "/proc/$FCPID/status"
kill $FCPID 2>/dev/null || true
# Compare with the qemu process for vss1: typically 40-80x larger.

# --- 8. Kata, if the host has it: one flag, a whole VM. ---------------
# Under Kubernetes this is runtimeClassName: kata in the Pod spec.
command -v kata-runtime >/dev/null && kata-runtime check || \
  echo "kata not installed -- see Appendix A"

# --- 9. Tear down. ----------------------------------------------------
cd ~ && rm -rf fc
\end{lstlisting}

\textbf{What to notice.} Three boot times spanning roughly two orders of
magnitude, and the middle one has the same isolation boundary as the slowest.
That single fact is why this chapter exists, and why Lambda is built the way it
is.
\end{lab}

\begin{checkpoint}
\begin{tightnum}
  \item Firecracker boots in 125 ms where a conventional guest takes 25 s. Name
        the three things that were removed, and say which of them a guest
        operating system must be willing to give up.
  \item A workload runs 18\% slower under gVisor than under runc, while another
        is unaffected. What distinguishes them, and what would you measure to
        confirm it?
  \item A cluster runs 200 internal services and accepts customer-supplied
        build jobs. Describe the runtime arrangement you would deploy and the
        single Kubernetes field that expresses it.
\end{tightnum}
\end{checkpoint}

\begin{chaptersummary}
\begin{tight}
  \item Isolation, density and start-up form a three-way trade. This chapter is
        the set of designs that refuse to pick only two.
  \item A microVM deletes the emulated 1990s PC. Firecracker is 50{,}000 lines,
        boots in 125 ms and costs under 5 MB --- the cost was compatibility,
        not virtualization.
  \item Firecracker is itself jailed in namespaces, a cgroup and a seccomp
        filter: containers containing a hypervisor. The layers compose.
  \item gVisor implements Linux in user space, cutting the host's exposed
        system calls from about 350 to about 60, at a system-call-rate-dependent
        performance cost and with imperfect coverage.
  \item Kata gives every Pod its own virtual machine behind the ordinary OCI
        interface, selectable per workload with RuntimeClass.
  \item Unikernels and Wasm go further still; the first has never paid down its
        operational cost, and the second is young.
  \item Isolation is now a per-workload setting rather than an architectural
        commitment. Most clusters should mix.
\end{tight}
\end{chaptersummary}

\begin{reviewq}
  \item \textbf{(MCQ)} Firecracker achieves its boot time principally by:
        (a)~a faster kernel; (b)~removing firmware, PCI enumeration and legacy
        device emulation; (c)~running without KVM; (d)~sharing the host
        kernel.
  \item \textbf{(MCQ)} gVisor reduces risk by: (a)~running each container in a
        virtual machine; (b)~implementing most of the Linux system-call
        interface in user space; (c)~encrypting guest memory; (d)~removing the
        container's network namespace.
  \item \textbf{(Short)} Explain why Firecracker ships a jailer, given that its
        own code base is only 50{,}000 lines.
  \item \textbf{(Short)} State the trade gVisor makes, and the workload
        characteristic that decides whether it is acceptable.
  \item \textbf{(Short)} Why have unikernels not displaced containers, despite
        better figures on size, start-up and attack surface?
  \item \textbf{(Applied)} Redo the worked example for functions with an
        average lifetime of 4 s rather than 400 ms. Give the start-up overhead
        for all three runtimes and say whether your recommendation changes, and
        why.
  \item \textbf{(Applied)} A university wants to let students submit code that
        is compiled and run automatically on a shared cluster. Set out the
        isolation design you would propose, what you would measure before
        going live, and what you would do if the measured overhead were twice
        what you expected.
\end{reviewq}
